\section{UPPAAL's Sequential Algorithm}\label{sec:background}

We recall the sequential zone-graph exploration as implemented in UPPAAL and
in the sequential path of \tool{}.  This section is deliberately concrete: our
goal is to exhibit the exact data-structure accesses that force serialisation,
so the reader can see immediately what the parallel approach changes.

\subsection{Symbolic states and zones}

A \emph{symbolic state} of a timed automaton with clock set
$\clocks = \{x_1,\ldots,x_n\}$ is a pair $(\ell, Z)$ where $\ell \in \Loc$
is a control location and $Z$ is a \emph{zone}: a convex polyhedron over
$\mathbb{R}_{\geq 0}^n$ defined by constraints $x_i - x_j \leq b_{ij}$.
Zones are stored as Difference Bound Matrices~\cite{Dill89}: an
$(n{+}1)\times(n{+}1)$ matrix over $\mathbb{Z}\cup\{+\infty\}$, kept in
canonical (shortest-path-closed) form so that inclusion and intersection run
in $O(n^2)$.

The successor function $\succfn(\ell, Z)$ maps one symbolic state to a
set of symbolic states.  For each edge $e = (\ell, g, R, \ell')$ leaving
$\ell$ it:
\begin{enumerate}
  \item intersects $Z$ with the guard $g$;
  \item resets the clocks in $R$ to zero;
  \item intersects with the target invariant $\Inv(\ell')$;
  \item applies time elapse (\textsc{delay}) if $\ell'$ is neither urgent nor
        committed;
  \item applies LU$^+$ extrapolation $\extraLU(Z, L, U)$ to widen the zone
        and guarantee termination~\cite{Behrmann06};
  \item discards the result if it is empty.
\end{enumerate}
Each of these steps touches only the local DBM and the (globally read-only)
automaton structure.  The passed list is \emph{never consulted}.

\subsection{The sequential exploration loop}

\begin{algorithm}[t]
\caption{Sequential zone-graph BFS (UPPAAL-style)}\label{alg:seq}
\begin{algorithmic}[1]
\Require TA $\mathcal{A}$,\ target $(\ell^\star, \phi)$
\State $Z_0 \leftarrow \mathit{init}(\mathcal{A})$;\enspace
       $\Passed[\Init] \leftarrow \{Z_0\}$;\enspace
       $\Wait \leftarrow \langle (\Init, Z_0) \rangle$
\While{$\Wait \neq \emptyset$}
  \State $(\ell, Z) \leftarrow \Wait.\mathbf{pop}()$
  \ForAll{$(\ell', Z') \in \succfn(\ell, Z)$}
    \If{$\mathit{isTarget}(\ell', Z')$} \Return \textbf{true} \EndIf
    \If{$\exists Z_p \in \Passed[\ell'] : Z' \subseteq Z_p$}
      \textbf{continue}
    \EndIf
    \State $\Passed[\ell'] \mathrel{+}= Z'$;\quad
           $\Wait.\mathbf{push}(\ell', Z')$
                                              \Comment{\textit{shared writes}}
  \EndFor
\EndWhile
\Return \textbf{false}
\end{algorithmic}
\end{algorithm}

\Cref{alg:seq} is the exact sequential algorithm (see \cref{fig:seq-loop} for a
data-flow view).  Three operations touch shared mutable state:
\begin{itemize}
  \item \textbf{Pop} from $\Wait$ (line~2): one state per iteration.
  \item \textbf{Read} from $\Passed[\ell']$ (line~4): the subsumption check.
  \item \textbf{Write} to $\Passed[\ell']$ and $\Wait$ (line~6): only when a
        state is genuinely new.
\end{itemize}
Because line~4 reads the same data that line~6 writes, no two iterations
can overlap: iteration $i+1$ might see a zone that iteration $i$ is about to
write.  This \emph{read-after-write dependency between consecutive iterations}
is the sole source of serialisation.

\subsection{What LU\texorpdfstring{$^+$}{+} extrapolation guarantees}

The LU$^+$ operator widens each zone so that clock values above the maximum
upper-bound constant $U_i$ are treated as ``large'' and values below
$-L_i$ are treated as ``small''.  Two important consequences follow:
\begin{itemize}
  \item \textbf{Finite quotient:} the number of distinct extrapolated zones at
        any location is bounded by a constant that depends only on the
        automaton's clock constants, not on the depth of exploration.
  \item \textbf{Small passed lists:} $|\Passed[\ell]|$ is therefore bounded
        above by this constant.  In practice it is rarely more than a few
        hundred entries even for large models.
\end{itemize}
The second point is crucial for the efficiency of our sequential merge phase:
the subsumption check for one candidate state costs
$O(n^2 \cdot |\Passed[\ell]|)$, which is $O(n^2)$ up to a model-dependent
constant.
